% !TeX root = ../main.tex
\section{Background and related work}
\label{sec:background}

\subsection{Training dialogue agents against an LLM judge}
Reinforcement learning from human feedback \citep{christiano2017deeprl, ouyang2022instructgpt}
replaces a hand-specified reward with a learned model of human preference. Direct Preference
Optimization \citep[DPO;][]{rafailov2023dpo} removes the separate reward model, optimizing a
policy directly from $(\text{prompt}, \text{chosen}, \text{rejected})$ triples. Group-Relative
Policy Optimization \citep[GRPO;][]{shao2024deepseekmath} keeps an online loop but replaces
PPO's learned value function \citep{schulman2017ppo} with the standardized advantage of each
completion within a sampled group.

In specialized domains without preference data, the loop is closed with a large model acting
as judge \citep{zheng2023judging}. For MI specifically, \citet{yosef2024assessing} established
the pattern this work builds on: simulate patients with an LLM, and score sessions with an LLM
that fills in validated MI questionnaires. The attraction is that the questionnaires are
instruments clinical researchers already use; the risk is that the judge becomes an
optimization target that was never designed to withstand optimization pressure
\citep{gao2023scaling, skalse2022defining}. Known failure modes of LLM judges --- sycophancy
\citep{sharma2024sycophancy, perez2023discovering}, length bias
\citep{dubois2024lengthcontrolled} and self-preference \citep{panickssery2024selfpreference}
--- all bear on that risk. This chapter is concerned with a different question, but the same
caution motivates our insistence on a second, decoupled grader in \S\ref{sec:gen3}.

\subsection{Search over conversational futures}
Where a reward is delayed, search is the classical remedy. Tree-structured exploration has been
combined with preference learning for reasoning tasks --- MCTS-guided step-level preference data
\citep{xie2024mctsdpo}, preference trees for coding and mathematics \citep{yuan2024eurus}, and
prompt-based multi-role search \citep{yu2023promptmcts}. Selecting best and worst from a sampled
pool to form preference pairs \citep{pace2024westofn} and generating preferences online from the
current policy \citep{guo2024oaif, yuan2024selfrewarding} are the score-based counterparts.
\citet{chen2025broaden} reduce reliance on direct simulation by planning in a semantic space of
conversations.

Those results come from domains with verifiable or near-verifiable outcomes. MI has neither: the
objective is subjective, the horizon is a whole session, and the only available referee is
another language model. Whether search over conversational futures pays off \emph{there} is
exactly what the three generations disagree about.

\subsection{Preference-Tree Optimization and the look-ahead lever}
\label{sec:pto}
PTO \citep{baruch2025pto} is the framework this chapter evaluates; DPO is the loss it uses.
One iteration proceeds as follows. The current therapist policy $\pi_n$ converses with a
simulated patient $P$. At each therapist turn along a growing trunk, $M$ candidate replies are
sampled from $\pi_n$. Each candidate is scored by an oracle $O$, the best is appended to the
trunk (so the choice conditions everything downstream), and a preference pair is emitted
wherever the best beats the worst by more than a filter threshold $\tau$. The collected pairs
train $\pi_{n+1}$ by DPO, and the loop repeats.

\paragraph{The lever.}
\emph{Look-ahead depth} $\K$ controls one thing only: how much conversation the oracle sees
before it scores a candidate. Given a prefix $c$ ending on a patient turn and a candidate
therapist reply $t$, look-ahead simulates $\K$ further alternating utterances under the current
policy,
\[
c \;+\; t \;+\; P(c{+}t) \;+\; \pi_n(\cdot) \;+\; P(\cdot) \;+\; \cdots
\]
and hands the \emph{extended} transcript to the oracle. At $\K{=}0$ the oracle scores $c{+}t$
alone. Nothing else changes: not the loss, not the branching factor, not the filter, not the
number of updates. The stated rationale is that scoring $c{+}t$ rewards replies that look good
in isolation, whereas scoring the extended trajectory rewards replies that \emph{lead
somewhere good} under the policy that will actually continue the conversation.

Because $\K$ touches only the scoring context, it is in principle meaningful for any optimizer
that consumes oracle scores. \textbf{This chapter nonetheless evaluates it entirely within
PTO.} That is the design choice that makes the three generations comparable: all three ran PTO,
so look-ahead is the single lever that varies across them. Whether look-ahead behaves the same
way under a different optimizer is a separate question, and one this chapter does not answer
(\S\ref{sec:limitations}).

\paragraph{What was claimed.}
The ICLR experiment ran $\K \in \{0,5\}$ for seven iterations each against a Llama-2-7B base and
reported that every PTO model beat the untrained baseline, that $\K{=}5$ models scored higher
than $\K{=}0$ models, and that the best $\K{=}5$ model had the lowest variance and the shortest
conversations \citep{baruch2025pto}. The comparison behind the $\K$ claim was a post-hoc Tukey
HSD between the best model of each arm. The paper was careful about its own evidence --- it
noted that the best $\K{=}5$ model's advantage over the best $\K{=}0$ model reached significance
only on the working-alliance questionnaire --- but the abstract and discussion generalized to
look-ahead being the component that produced ``the most stable and high-scoring results''. It is
that generalization the present chapter tests.
